%%%%%%%%%%%%%%%%%%%%%%%%%%%%%%%%%
%
%       EXERCÍCIO
%
%%%%%%%%%%%%%%%%%%%%%%%%%%%%%%%%%

\definevalorquestao

\begin{exercicioBanco}[\valorquestao]
No mercado de segunda mão, considera-se que certo \textit{smartphone} perde \(\dfrac{1}{4}\) do seu valor em relação ao ano anterior a cada ano de uso. Após 4 anos, o proprietário de um aparelho cujo preço original era R\$ 5.120,00 deseja trocá-lo por um modelo novo e avalia duas opções: vendê-lo no mercado de segunda mão ou entregá-lo na loja onde comprará o novo modelo, que o recebe por \(35\%\) do seu valor original de compra.

Analisando financeiramente as opções, a alternativa mais vantajosa para esse consumidor é:

% Define as alternativas
\newcommand{\alternativas}{%
    \begin{center}
        \begin{tabularx}{\textwidth}{X}
            (a) Entregar o aparelho na loja, pois obterá R\$ 172,00 a mais do que no mercado de segunda mão. \\[3pt]
            (b) Vender no mercado de segunda mão, pois obterá R\$ 172,00 a mais do que na oferta da loja. \\[3pt]
            (c) Entregar o aparelho na loja, pois obterá R\$ 292,00 a mais do que no mercado de segunda mão. \\[3pt]
            (d) Vender no mercado de segunda mão, pois obterá R\$ 92,00 a mais do que na oferta da loja. \\[3pt]
            (e) Ambas as opções resultam exatamente no mesmo valor financeiro.
        \end{tabularx}
    \end{center}
}

% Define a resposta correta
\newcommand{\resposta}{A}

% Lógica condicional para exibição
\ifdefstring{\atividade}{lista}{%
    \alternativas
    \vspace{0.5em}
    
    \noindent\textbf{Resposta:} letra \textbf{\resposta}.
}{%
    \ifdefstring{\modo}{objetivo}{%
        \alternativas
    }{%
        % modo = discursiva → não mostra alternativas
    }
}
\end{exercicioBanco}

